\section{Conclusiones}
\label{sec:conclusiones}

Los resultados permiten responder directamente la pregunta planteada:

\begin{enumerate}

    \item Mejorar termoestabilidad no obligó a sacrificar eficiencia. En los
    cuatro casos se obtuvo \(I_S>1\) e \(I_C>1\).

    \item Bajo la condición base común, ENG generó entre \(4.58\times\) y
    \(15.29\times\) la producción de WT. La magnitud depende del cuello de
    botella de cada aplicación.

    \item Cuando cada variante se llevó a su mejor punto dentro del rango
    evaluado, ENG mantuvo una ventaja de \(2.78\times\) a \(5.95\times\).
    Por tanto, la ventaja no depende únicamente de haber elegido una condición
    desfavorable para WT.

    \item El punto de máxima actividad no siempre coincide con el de máxima
    producción. Biomasa lo muestra con claridad: ENG presenta máxima actividad
    alrededor de \(75\,^{\circ}\mathrm{C}\), pero su mayor producto acumulado
    aparece alrededor de \(70\,^{\circ}\mathrm{C}\).

    \item La decisión industrial debe atacar el cuello de botella concreto:
    vida operativa en alimentos, estabilidad y tolerancia en bagazo, ventana
    térmica en biomasa y selectividad en glucósidos.
\end{enumerate}

La conclusión general es:

\begin{equation}
\text{Ingeniería Enzimática}
\longrightarrow
\text{Mejora Molecular}
\longrightarrow
\text{Control del Cuello de Botella}
\longrightarrow
\text{Mayor Producción Útil}.
\end{equation}

La ingeniería enzimática puede mejorar simultáneamente estabilidad y eficiencia,
pero su valor industrial depende de operar la variante en condiciones que
conviertan esa mejora en producto
\citep{siddiqui_defying_2017,hou_tradeoff_2023}.
